% GENERATED FILE. Edit canonical lessons, then run npm run generate:books.
% canonical-source-hash: fnv1a64:e6e28200eba8eafb
% canonical-lessons: SA-C145-nandati, SA-C145-krandati, SA-C145-garjati, SA-C145-kujati, SA-C145-jalpati

\chapter{Is glad, Cries out, Roars, Coos, Chatters}
\label{ch:sa-is-glad-cries-out-roars-coos-chatters}
\hlchaptermodality{145}

\begin{chapteropening}
\textbf{By the end of this chapter:} \emph{I can say is glad, cries out, roars, coos and chatters.}

Complete the last lesson of chapter 145: I can say is glad, cries out, roars, coos and chatters.
\end{chapteropening}

\section[\texorpdfstring{\sk{नन्दति} (nandati)}{नन्दति (nandati)}]{\sk{नन्दति} (nandati) — is glad}

\label{lesson:SA-C145-nandati}



\begin{quote}
Before the new one: say the Sanskrit for puts in, then the Sanskrit for joins.
\end{quote}

\subsection*{You'll want to know: \sk{नन्दति}}
\textbf{\sk{नन्दति}} — \emph{nandati} — \textquotedblleft{}is glad\textquotedblright{}.

\begin{grammarlens}[title={What you've built}]
One more word for this chapter.
\end{grammarlens}

\subsection*{Your turn}
\begin{itemize}
\raggedright
  \item \emph{Say it:} \emph{nandati}
  \item \emph{Say it:} \emph{nandati}, once more
  \item \emph{Say it:} say \emph{nandati}
  \item \emph{Recall:} say the Sanskrit for puts in, then the Sanskrit for joins, then say \emph{nandati} again
\end{itemize}

\begin{tcolorbox}[breakable,colback=teal!4,colframe=teal!35!black,title={Before you move on}]
What does \sk{नन्दति} mean? (\textquotedblleft{}Is glad\textquotedblright{}.) Say it once more.
\end{tcolorbox}
\section[\texorpdfstring{\sk{क्रन्दति} (krandati)}{क्रन्दति (krandati)}]{\sk{क्रन्दति} (krandati) — cries out, wails}

\label{lesson:SA-C145-krandati}



\begin{quote}
Before the new one: say the Sanskrit for joins, then the Sanskrit for is glad.
\end{quote}

\subsection*{You'll want to know: \sk{क्रन्दति}}
\textbf{\sk{क्रन्दति}} — \emph{krandati} — \textquotedblleft{}cries out, wails\textquotedblright{}.

\begin{grammarlens}[title={What you've built}]
One more word for this chapter.
\end{grammarlens}

\subsection*{Your turn}
\begin{itemize}
\raggedright
  \item \emph{Say it:} \emph{krandati}
  \item \emph{Say it:} \emph{krandati}, once more
  \item \emph{Say it:} say \emph{krandati}
  \item \emph{Recall:} say the Sanskrit for joins, then the Sanskrit for is glad, then say \emph{krandati} again
\end{itemize}

\begin{tcolorbox}[breakable,colback=teal!4,colframe=teal!35!black,title={Before you move on}]
What does \sk{क्रन्दति} mean? (\textquotedblleft{}Cries out, wails\textquotedblright{}.) Say it once more.
\end{tcolorbox}
\section[\texorpdfstring{\sk{गर्जति} (garjati)}{गर्जति (garjati)}]{\sk{गर्जति} (garjati) — roars, thunders}

\label{lesson:SA-C145-garjati}



\begin{quote}
Before the new one: say the Sanskrit for is glad, then the Sanskrit for cries out.
\end{quote}

\subsection*{You'll want to know: \sk{गर्जति}}
\textbf{\sk{गर्जति}} — \emph{garjati} — \textquotedblleft{}roars, thunders\textquotedblright{}.

\begin{grammarlens}[title={What you've built}]
One more word for this chapter.
\end{grammarlens}

\subsection*{Your turn}
\begin{itemize}
\raggedright
  \item \emph{Say it:} \emph{garjati}
  \item \emph{Say it:} \emph{garjati}, once more
  \item \emph{Say it:} say \emph{garjati}
  \item \emph{Recall:} say the Sanskrit for is glad, then the Sanskrit for cries out, then say \emph{garjati} again
\end{itemize}

\begin{tcolorbox}[breakable,colback=teal!4,colframe=teal!35!black,title={Before you move on}]
What does \sk{गर्जति} mean? (\textquotedblleft{}Roars, thunders\textquotedblright{}.) Say it once more.
\end{tcolorbox}
\section[\texorpdfstring{\sk{कूजति} (kūjati)}{कूजति (kūjati)}]{\sk{कूजति} (kūjati) — coos, warbles}

\label{lesson:SA-C145-kujati}



\begin{quote}
Before the new one: say the Sanskrit for cries out, then the Sanskrit for roars.
\end{quote}

\subsection*{You'll want to know: \sk{कूजति}}
\textbf{\sk{कूजति}} — \emph{kūjati} — \textquotedblleft{}coos, warbles\textquotedblright{}.

\begin{grammarlens}[title={What you've built}]
One more word for this chapter.
\end{grammarlens}

\subsection*{Your turn}
\begin{itemize}
\raggedright
  \item \emph{Say it:} \emph{kūjati}
  \item \emph{Say it:} \emph{kūjati}, once more
  \item \emph{Say it:} say \emph{kūjati}
  \item \emph{Recall:} say the Sanskrit for cries out, then the Sanskrit for roars, then say \emph{kūjati} again
\end{itemize}

\begin{tcolorbox}[breakable,colback=teal!4,colframe=teal!35!black,title={Before you move on}]
What does \sk{कूजति} mean? (\textquotedblleft{}Coos, warbles\textquotedblright{}.) Say it once more.
\end{tcolorbox}
\section[\texorpdfstring{\sk{जल्पति} (jalpati)}{जल्पति (jalpati)}]{\sk{जल्पति} (jalpati) — chatters}

\label{lesson:SA-C145-jalpati}



\begin{quote}
Before the new one: say the Sanskrit for roars, then the Sanskrit for coos.
\end{quote}

\subsection*{You'll want to know: \sk{जल्पति}}
\textbf{\sk{जल्पति}} — \emph{jalpati} — \textquotedblleft{}chatters\textquotedblright{}.

\begin{grammarlens}[title={What you've built}]
One more word for this chapter.
\end{grammarlens}

\subsection*{Your turn}
\begin{itemize}
\raggedright
  \item \emph{Say it:} \emph{jalpati}
  \item \emph{Say it:} \emph{jalpati}, once more
  \item \emph{Say it:} say \emph{jalpati}
  \item \emph{Recall:} say the Sanskrit for roars, then the Sanskrit for coos, then say \emph{jalpati} again
\end{itemize}

\begin{tcolorbox}[breakable,colback=teal!4,colframe=teal!35!black,title={Before you move on}]
What does \sk{जल्पति} mean? (\textquotedblleft{}Chatters\textquotedblright{}.) Say it once more.
\end{tcolorbox}
